% 【综合论述】所属：第9章（由 chapters/revised/09-comparative.tex \input 引入）
\minitopic{双向批评要求双方都可能改变}若比较只问古代文本是否符合现代知识分析，现代框架始终不承担风险；若只用传统整体智慧批评分析哲学狭窄，传统概念也无需接受精确反例。双向批评先让每方以最强版本提出问题，再要求对方暴露边界。命题知识分析可以迫使厚概念区分真值、能力和责任；修身与实践传统则可以迫使命题分析说明知识为何与注意、言说角色和行动相连。比较成功的标志不是达成同一词汇，而是双方问题都变得更清楚。

\minitopic{不可通约是一项发现而非禁令}两种框架可能缺少无损共同尺度，例如一边评价孤立命题状态，另一边评价长期工夫和人格转变。强行统一会歪曲材料，但宣布完全不可比较也过快。研究者仍可比较错误如何被发现、学习怎样改变主体、共同体如何确认成就，以及各框架无法表达什么。不可通约报告应指出具体损失发生在哪里，并说明这一边界怎样限制结论；它不是“各说各话”的终点。

\minitopic{规范建构必须由研究者署名}把多个传统资源结合成新理论可以富有创造性，例如用反运气条件约束命题真值联系，用修身资源解释注意和品格，再用社会认识论分析制度依赖。但这套组合不是任何一位历史作者已经完整主张的理论。研究者应明确选择理由、内部张力和反例责任。规范建构越成功，越不需要借“古人早已说过”取得权威；它应以当前论证、证据和实践效果接受评价。

\minitopic{九个问题构成一张可迁移地图}本书从知识分析、怀疑与运气、理由结构、认识来源、社会依赖、形式模型、认识价值、数字环境到比较方法，提供九组互相连接的问题。面对一个新案例，不必机械套用全部术语，而应判断失败发生在哪一层：命题是否为真，理由是否实际支撑信念，方法是否抵抗近邻错误，来源链是否可追踪，群体是否保存异议，数字模型是否稳健，认识成果服务何种价值，以及采用的概念是否跨语境失真。这张地图帮助定位，不是自动评分器。

\begin{argumentbox}
形成个人立场时建立“理论档案”：核心主张；适用范围；最强正例；最强反例；为回应反例付出的代价；来自其他传统的修正压力；可能改变立场的新证据。档案使立场成为可修订研究方案，而不是一组不可动摇的主义标签。
\end{argumentbox}

\minitopic{比较也要接受经验与语言证据}不同群体对案例的判断、翻译选择和实践背景可能变化。经验差异不会直接决定规范真理，却能削弱“人人都这样使用概念”的描述主张，也能揭示问卷预设了哪种主体模型。研究者应区分语言使用、直觉分布、历史解释与规范理由，使用各自证据。跨文化数据若只是把英语案例直译并比较票数，很可能测到理解难度和教育背景；高质量研究需要共同设计案例、检查概念功能并报告不可比较部分。

\minitopic{认识谦逊不等于停止判断}承认翻译有损失、证据有限和理论可能修正，不意味着所有解释同样好。研究者仍可根据文本覆盖、论证一致、解释范围和反例承受力作出有程度的判断。谦逊表现为让主张强度与证据相称、保存替代解释、说明何种材料会改变结论。只说“一切都有可能”逃避了评价，也不能指导后续研究。

\minitopic{问题驱动不是取消经典文本}以问题组织教材，能够把原典放入真实论证，而不是按年代排成名单。读者先知道某段文本在回答什么，才能判断其前提、反例和当代意义；同时，文本阻力又会迫使既有问题改写。问题与文本因此形成循环：问题引导阅读，内部解释限制比较，比较产生新问题。只要转换被公开标明，这种循环不是任意，而是理解逐步加深的机制。

\minitopic{全书结论}认识论不只问“我有没有知识”，还问成功为何不偶然、理由怎样支撑信念、来源和制度如何工作、概率如何表达不确定、理解与智慧为何重要、技术怎样重塑证据，以及不同传统如何重新安排这些问题。没有一种短定义能够独自完成全部任务。更可靠的目标是一套可审计、可比较、可修订的探究：明确命题，追踪理由，检验反例，保存来源，报告不确定，分配责任，并让新证据真正能够改变结论。
